\section{Lifting every marked arc in row B}

% BEGIN B82R1
\paragraph{B82R1. From a target cut vector to an integral numerator.}
At a cut of a marked row B word let $W_\nu=\Lambda^{-1}
A\varphi^\ell N_{\nu,c-\ell}$, where $N_{\nu,c}$ is the
invariant Wach lattice of the source ($\nu=1$) or target
($\nu=2$) and $A=O[[v]]$. Write $\mathcal L_{\nu,s}$ and
$\mathcal B_{\nu,s}$ for the true and rounded Taylor lattices
of D82P1--D82P3. Let $D_{\nu,s}$ denote the constant periods
at level $s$. Inclusions at a level refer to all its primitive
branches. Take a target cut vector $w\in W_2$ for which
$(\gamma-1)w\in vN_{2,c}$, as in H82A3.

The generator preserves $W_\nu$ and acts trivially on
$W_\nu/vW_\nu$. Indeed in its basis
$\varphi^\ell(e)/\Lambda$ the action on the numerator is
congruent to the identity modulo $\varphi^\ell(v)$ by J82F1;
this parameter is divisible by $v$. The action on each
principal factor $(P_s)$ preserves its ideal: the change of
variable permutes the primitive zeros, and division by the
associated distinguished polynomial gives divisibility in both
directions. The unit $\gamma(P_s)/P_s$ has constant term one
since $P_s(0)=p$. The same holds for the denominator ratio.
The actions on coefficients and on $N_{\nu,c}$ are also the
identity modulo $v$ by the Wach condition. In particular
$vN_{\nu,c}$ is stable, and polynomials $Q$ over $O$ satisfy
\[
           Q(\gamma)w\equiv Q(1)w\pmod{vN_{2,c}} .
\]

Take the target polynomial $P$ of K82P2. Localizing the
difference of $w$ gives an element of $\mathcal L_{2,s}$,
by D82P1 and the compatibility of the generator action with
the primitive expansion. K82P1 then puts the jet of $w$ in
$D_{2,s}+\mathcal L_{2,s}$. By D82P3 this lies in
$D_{2,s}+P_s^{-g}\mathcal B_{2,s}$. Consequently $x=P(\gamma)w$
lies in $W_2$ and its jets belong to
$D_{2,s}+\mathcal B_{2,s}$ by K82P2. Its residue modulo
$vN_{2,c}$ is $P(1)w$, so $(\gamma-1)x\in vN_{2,c}$ too.

Use the target direction of J82F2 to choose one integral model
numerator $y$ whose jets, after division by $\Lambda$ in the
constant frame, agree with $p^\sigma x$ up to the error space
at every level. D82P3 puts that error in the rounded lattice.
The jet characterization in A82N1 now gives
$y\in\mathcal I_2$. The constants permitted there are common
to the branches of each level by K82P1.
% END B82R1

% BEGIN B82R2
\paragraph{B82R2. An exact physical preimage after the correction in the Wach lattice.}
Apply A82N5 and P82U3 to this numerator. With $f_0$ the
map in the constant frames there are $y_1\in\mathcal I_1$
and $\nu\in\mathcal N_2$ such that
\[
        p^{2a_0}y=f_0y_1+p^{2a_0}\nu .
\]
By J82F2 choose $x_1\in W_1$ with jets agreeing with those
of $p^\sigma y_1$ in the constant model up to the frame errors
(the model divides a numerator by $\Lambda$). Its jets belong
to $D_{1,s}+\mathcal B_{1,s}$. Put
\[
             \delta=p^{2\sigma+2a_0}x-f(x_1)\in W_2 .
\]
The map preserves the Taylor orders of the frame errors: in
the two period bases it is the constant linear map of J82F3,
and the same scalar $\Lambda$ is used on both sides. The
displayed equation and both partner comparisons show that the
jets of $\delta$ belong to $\mathcal B_{2,s}$. Indeed they are
the jets of $p^{\sigma+2a_0}\nu$ plus frame errors in that
lattice. Writing $\Theta=\prod_{s<\ell}P_s^g$, D82P3 and
D82P2 therefore give $\Theta\delta\in N_{2,c}$.

By R82W4 lift $p^R\Theta\delta$ to $n_1\in N_{1,c}$.
Each $P_s$ is invertible in the integral \'etale coefficient
ring: modulo the coefficient uniformizer it is a unit times
a power of the invertible variable, and inversion lifts by
completeness. Thus $\delta_1=n_1/\Theta$ exists there and
belongs to $W_1$ by D82P2. Its jets lie in
$P_s^{-g}\mathcal L_{1,s}$ by the same test. The vector
\[
          z=p^R x_1+\delta_1\in W_1
       \quad\hbox{satisfies}\quad
          f(z)=p^{2\sigma+2a_0+R}x .
\]
Its jets belong to $D_{1,s}+P_s^{-g}\mathcal L_{1,s}$.
Here the cost $R$ occurs once, on the Wach correction to a
single source numerator valid at every primitive level.
% END B82R2

% BEGIN B82R3
\paragraph{B82R3. The source cut congruence and the whole guarded row B.}
Apply the source polynomial $P'$ of K82P3 to $z$ and put
$w'=P'(\gamma)z\in W_1$. Its jets belong to
$D_{1,s}+\mathcal L_{1,s}$. The generator fixes these
constants and preserves the true Taylor lattices, by K82P1.
Hence $(\gamma-1)w'$ passes the true lattice test at every
level. It also belongs to $vW_1$ by B82R1. The two tests
together give $(\gamma-1)w'\in vN_{1,c}$ by D82P2. Thus
$w'$ is an actual source cut vector. Functoriality and the
residue calculation of B82R1 give
\[
 f(w')\equiv
     p^{2\sigma+2a_0+R}P'(1)P(1)w\pmod{vN_{2,c}}.
\]
K82P4 writes the product as a unit times $p^k$ for an
integer $0\le k\le2\kappa$. This proves the congruence
of H82A3 for every target cut vector on each marked row B
orbit. That argument now produces a finite source lift at
$e_{\rm par}$ for every finite target class, with the
coefficient and integral lattice of the tested surjection.

Choose the free row constant $C_2\ge122$ before applying
F78D2 and choosing the common family cutoff. I82B2--I82B3
bound the resulting normalized length at every marked orbit
by the rational row bound stored for this choice. I82B1 gives
it at every unmarked orbit. The choices of coefficient
extension, surjection, place and good point were arbitrary.
G82I1--G82I2 therefore give the original middle-range
inequality on this hub at \emph{every} guarded inert row B
index. The length estimate on marked orbits uses the actual
cut cocycles, physical frames and integral inverse images on
the same tested map.
% END B82R3
